% Pursuing Stacks (cotes 134-2 to 134-8): archive cote 134-7, site page 20 (PDF page 21), typescript page 462.
% typescript. Draft compared: ps6.tex:930-988.
% AI-transcribed by Claude Opus 5.5 (claude-opus-5-5, adaptive thinking, effort high), Message Batches API, batch ps3 (msgbatch_01XdFmuVVBAcgc7EFxdtCqsj), 2026-11-12; not yet checked by eye.
\scanpage{134-7/20}
\mspage{462}

affine groups over a field (here the field of fractions $\underline{Q}$ of $\underline{Z}$ ) tells
us that $G_{\underline{Q}}$ is a \uline{nilpotent} algebraic group,  and that therefore it admits
a composition series with factors isomorphic to the additive group $G_{a\,\underline{Q}}$ ;
presumably the same dévissage then can be obtained over the base $\underline{Z}$ ,  and
using induction on the length of the composition series,  and the Hoch-%
schild-Serre type of relations (traditionally expressed by a spectral
sequence)  between group cohomology of a group,  quotient group and corres
ponding subgroup,  we should get the wished for quasi-isomorphisms (w).

Take now however the simplest case of a monoid which is'nt a group,
namely the multiplicative law on the affine line,  given by the polynomial
law
\[ W(\underline{Z}) \text{ x } W(\underline{Z}) \hand{\longrightarrow} W(\underline{Z}) \quad : \quad (x,y) \hand{\longmapsto} xy. \]
\margin{\textcircled{$\times$}} the corresponding ``discrete'' monoid is just $\underline{Z}^{\hand{(x)}}$ ,  namely the integers
with multiplication,  it's $H^{1}$ with coefficients in $\underline{Z}$  is just the group
of all homomorphisms
\[ (*) \qquad \underline{Z}^{\hand{(x)}} \hand{\to} \underline{Z}^{\hand{(+)}} = \underline{Z} \quad , \]

and denoting by $\underline{P}$ the set of all primes and using the prime decomposition
of integers,  we find that
\[ \mathrm{Hom}(\underline{Z}^{\hand{(x)}},\underline{Z}^{\hand{(+)}}) \hand{\simeq} \underline{Z}^{P} \quad , \]
to a family $(n_{p})_{p \in P}$ of integers being associated the homomorphism
\[ \hand{\pm\prod_{p\in P}} p^{\hand{\alpha}_{p}} \hand{\mapsto \sum_{p\in P} \alpha}_{p}\, n_{p} \quad . \]
On the other hand,  the schematic $H^{1}$ consists \over{\illegible{1}}{of} all  homomorphisms (*)
which can be expressed by a polynomial,  hence induce a homomorphism of
algebraic group schemes  $\underline{G}_{m} \hand{\to} \underline{G}_{a}$ ,  and it is well known (and immediately
checked) that there is only the zero homomorphism !

Thus,  it turns out that the assumptions made yesterday on $X_{*}$ ,
in order for the ``linearization theorem'' (!) to hold\hand{,} namely the maps (w)
to be quasi-isomorphisms,  are definitely not strong enough yet\over{.}{\unsure{!}} One
may think of throwing in the extra condition  $X_{1}=e$ ,  so as to rule out
monoids\inserted{altogether,} (and even groups, too bad !), \deleted{altogether,}  but I do'nt think this
helps at all (did'nt try though to make a counterexample\hand{).}\deleted{\xout{still}).}  On the
other hand,  just restrict\inserted{ing}\deleted{\unsure{to}} Kan-Postnikov complexes seems rather awk-%
ward,   we defini\xout{\illegible{1}}tely do'nt want to drag along P\over{\illegible{1}}{o}stnikiv fibrations as
a compulsory ingredient of the complexes we work with.  The idea which
comes up here is just to \hand{``}drop P\over{\illegible{1}}{o}stnikow and keep Kan\hand{''} -  namely \uline{introduce}
\uline{a Kan type condition on semisimplicial complexes in U(k)}.  If we mimick

\margin{20}
